\chapter{Chapter09 — Hypercontractivity Exercises}
%%% generated by the blueprint pipeline from FABL.Chapter09.HypercontractivityExercises
\section{Overview}

Declarations of the Lean module \texttt{FABL.Chapter09.HypercontractivityExercises}.

\section{Declarations}

\begin{definition}[Degree truncation]
\lean{FABL.degreeTruncation}
\statementsource{Analysis-of-Boolean-Functions-by-Ryan-ODonnell}{
  pdf-ffbe9c7b3306-p015-b014,
  pdf-ffbe9c7b3306-p028-b002,
  pdf-ffbe9c7b3306-p102-b005,
  pdf-ffbe9c7b3306-p271-b007
}
\label{FABL.degreeTruncation}
\leanok
\uses{FABL.fourierCoeff, FABL.monomial}
For $k\in\mathbb{N}$ and a function $f:\{-1,1\}^n\to\mathbb{R}$, define its degree-$k$
truncation by $\sum_{S\subseteq[n],\,|S|\le k}\widehat f(S)\,\chi_S(x)$, where $\widehat
f(S)$ is the Fourier coefficient of $f$ at $S$ and $\chi_S$ is the monomial indexed by
$S$.
\end{definition}

\begin{theorem}[Fourier coefficients of degree truncation]
\lean{FABL.fourierCoeff_degreeTruncation}
\label{FABL.fourierCoeff_degreeTruncation}
\leanok
\uses{FABL.degreeTruncation, FABL.expect_monomial_mul, FABL.fourierCoeff, FABL.monomial}
\difficulty{4}
Let $n,k\ge 0$, let $f$ be a real-valued function on the $n$-dimensional sign cube, and
let $T$ be a set of coordinates. The Fourier coefficient at $T$ of the degree-$k$
truncation of $f$ equals the Fourier coefficient of $f$ at $T$ when $|T|\le k$, and
equals $0$ when $|T|>k$.
\end{theorem}

\begin{theorem}[Degree of Fourier truncation]
\lean{FABL.fourierDegree_degreeTruncation_le}
\label{FABL.fourierDegree_degreeTruncation_le}
\leanok
\uses{FABL.degreeTruncation, FABL.fourierCoeff_degreeTruncation, FABL.fourierDegree, FABL.fourierDegree_le_iff}
\difficulty{3}
For a natural number $n$, a natural number $k$, and a function
$f\colon\{-1,1\}^n\to\mathbb{R}$, form the sum of the Fourier terms of $f$ indexed by
subsets of size at most $k$. The resulting function has Fourier degree at most $k$.
\end{theorem}

\begin{theorem}[Self-adjointness of degree truncation]
\lean{FABL.degreeTruncation_selfAdjoint}
\label{FABL.degreeTruncation_selfAdjoint}
\leanok
\uses{FABL.degreeTruncation, FABL.fourierCoeff_degreeTruncation, FABL.plancherel, FABL.uniformInner}
\difficulty{4}
Let $n,k\in\mathbb{N}$ and let $f,g:\{-1,1\}^n\to\mathbb{R}$. The uniform inner product
of the degree-$k$ truncation of $f$ with $g$ equals the uniform inner product of $f$
with the degree-$k$ truncation of $g$.
\end{theorem}

\begin{theorem}[Uniform lp norm two degree truncation le]
\lean{FABL.uniformLpNorm_two_degreeTruncation_le}
\statementsource{Analysis-of-Boolean-Functions-by-Ryan-ODonnell}{pdf-ffbe9c7b3306-p271-b007}
\label{FABL.uniformLpNorm_two_degreeTruncation_le}
\leanok
\uses{FABL.degreeTruncation, FABL.fourierCoeff_degreeTruncation, FABL.parseval, FABL.uniformLpNorm, FABL.uniformLpNorm_nonneg, FABL.uniformLpNorm_two_sq_eq_uniformInner}
\difficulty{4}
% dataset-pipeline: stub -- prose pending build_informal_statements.py
Statement pending regeneration.
\end{theorem}

\begin{theorem}[Norm bound for low-degree truncation]
\lean{FABL.degreeTruncation_norm_le_of_two_le}
\proofsource{Analysis-of-Boolean-Functions-by-Ryan-ODonnell}{pdf-ffbe9c7b3306-p271-b007}
\label{FABL.degreeTruncation_norm_le_of_two_le}
\leanok
\uses{FABL.degreeTruncation, FABL.fourierDegree_degreeTruncation_le, FABL.lowDegree_uniformLpNorm_le, FABL.uniformLpNorm, FABL.uniformLpNorm_mono, FABL.uniformLpNorm_two_degreeTruncation_le}
\difficulty{4}
Let $f$ be a real-valued function on the $n$-dimensional sign cube, let $k$ be a natural
number, and let $q$ be a real number with $2\le q$. The uniform $L^q$ norm of the
Fourier truncation of $f$ to degrees at most $k$ is at most $(\sqrt{q-1})^k$ times the
uniform $L^q$ norm of $f$.
\end{theorem}

\begin{theorem}[Norm bound for degree truncation]
\lean{FABL.degreeTruncation_norm_le_two_of_two_le}
\label{FABL.degreeTruncation_norm_le_two_of_two_le}
\leanok
\uses{FABL.degreeTruncation, FABL.fourierDegree_degreeTruncation_le, FABL.lowDegree_uniformLpNorm_le, FABL.uniformLpNorm, FABL.uniformLpNorm_two_degreeTruncation_le}
\difficulty{3}
Let $n,k\in\mathbb{N}$, let $f:\{-1,1\}^n\to\mathbb{R}$, and let $q\in\mathbb{R}$
satisfy $2\le q$. The uniform $L^q$ norm of the degree-$k$ truncation of the Fourier
expansion of $f$ is at most $(\sqrt{q-1})^k$ times the uniform $L^2$ norm of $f$.
\end{theorem}

\begin{theorem}[Norm bound for low-degree truncation]
\lean{FABL.degreeTruncation_norm_le_of_one_lt_le_two}
\label{FABL.degreeTruncation_norm_le_of_one_lt_le_two}
\leanok
\uses{FABL.degreeTruncation, FABL.degreeTruncation_norm_le_of_two_le, FABL.degreeTruncation_selfAdjoint, FABL.selfAdjoint_dual_bound, FABL.uniformLpNorm}
\difficulty{4}
Let $f$ be a real-valued function on the $n$-dimensional sign cube, let $k$ be a natural
number, and let $q$ satisfy $1<q\le 2$. The uniform $L^q$ norm of the Fourier truncation
of $f$ to degrees at most $k$ is at most $(1/\sqrt{q-1})^k$ times the uniform $L^q$ norm
of $f$.
\end{theorem}

\begin{theorem}[Fourier truncation recovers the function]
\lean{FABL.degreeTruncation_eq_self_of_fourierDegree_le}
\label{FABL.degreeTruncation_eq_self_of_fourierDegree_le}
\leanok
\uses{FABL.degreeTruncation, FABL.fourierCoeff, FABL.fourierCoeff_degreeTruncation, FABL.fourierDegree, FABL.fourierDegree_le_iff, FABL.fourier_expansion}
\difficulty{4}
Let $f$ be a real-valued function on the $n$-dimensional sign cube, and let $k$ be a
natural number. If the largest size of a set with a nonzero Fourier coefficient of $f$
is at most $k$, then truncating its Fourier expansion to sets of size at most $k$ gives
$f$ itself.
\end{theorem}

\begin{theorem}[Low-degree norm comparison]
\lean{FABL.lowDegree_uniformLpNorm_two_le_weak}
\label{FABL.lowDegree_uniformLpNorm_two_le_weak}
\leanok
\uses{FABL.degreeTruncation, FABL.degreeTruncation_eq_self_of_fourierDegree_le, FABL.degreeTruncation_norm_le_two_of_two_le, FABL.degreeTruncation_selfAdjoint, FABL.fourierDegree, FABL.selfAdjoint_dual_bound, FABL.uniformLpNorm}
\difficulty{4}
Let $f$ be a real-valued function on the $n$-dimensional sign cube whose Fourier degree
is at most $k$. If $1<p<2$, then
\[\|f\|_2\le \left(\frac{1}{\sqrt{p-1}}\right)^k\|f\|_p,\]
where the norms are uniform norms on the cube.
\end{theorem}

\begin{theorem}[Logarithmic inequality on the unit interval]
\lean{FABL.two_mul_one_sub_div_one_add_lt_neg_log}
\label{FABL.two_mul_one_sub_div_one_add_lt_neg_log}
\leanok
\difficulty{5}
For every real number $x$ with $0<x<1$, $\frac{2(1-x)}{1+x}<-\log x$.
\end{theorem}

\begin{theorem}[Exponential and reciprocal square-root bound]
\lean{FABL.exp_two_div_sub_one_lt_inv_sqrt_sub_one}
\label{FABL.exp_two_div_sub_one_lt_inv_sqrt_sub_one}
\leanok
\uses{FABL.two_mul_one_sub_div_one_add_lt_neg_log}
\difficulty{5}
For every real $p$ with $1<p<2$, one has $\exp(2/p-1)<1/\sqrt{p-1}$.
\end{theorem}
